\section{A two-feature tanh obstruction to lifting the scalar argument}
\label{sec:tanh}

We now return to tanh in \emph{every} hidden layer, an appended constant
input, and the prior \eqref{eq:prior}. The following is an energy
comparison, not a statement about the predictions of the wider posterior.
It shows why the scalar affine minimum cannot simply be inserted as
the minimum energy of the wider architecture.

\begin{proposition}[A cheap classifier consistent with the entire suffix]
\label{prop:tanh-shortcut}
Let $k\ge128$, $d=2k+1$, and $L=d+1=2k+2$.
For every width $N\ge2$, the Boolean function
\[
 g(x,y,z)=+1\quad\Longleftrightarrow\quad z=-1\text{ and }x_0=y_0
\]
has a strict-global-margin tanh implementation with every coefficient
bounded by two and
\[
 \kappa_{d,N}\le L+89+12\log d<2L.
\]
It fits every calibration example and every negative query $q_t$ from
Section~\ref{sec:selector}.
In contrast, every scalar tanh implementation of $\sgn A_t$, for
$k\le t\le M$, with weak global output margin $1/2$ has energy at least
$2L$.
\end{proposition}
\begin{proof}
Every positive calibration point is diagonal. At every negative carry
point, $x_0\ne y_0$: either bit zero is the highest differing bit, or
it is one of the lower carry bits. Every $q_t$ has $z=+1$.
Thus $g$ fits the entire stated sequence.

Put
\[
 \lambda=(82d)^{-1/2},\qquad
 s=x_0-y_0+\frac72+\frac52z,\qquad
 r=\frac{\tanh(4\lambda)}{\tanh(2\lambda)}
   =\frac2{1+\tanh^2(2\lambda)}.
\]
Consider the two-feature combination
\[
 F(s)=\tanh(2\lambda s)-r\tanh(\lambda s).
\]
Writing $a=\tanh(\lambda s)$ and $b=\tanh(2\lambda)$, the double-angle
identity gives
\begin{equation}\label{eq:tanh-two}
 F(s)=\frac{2a(b^2-a^2)}{(1+a^2)(1+b^2)}.
\end{equation}
It is odd, positive for $0<s<2$, and negative for $s>2$.
The possible $s$ values are $-1,1,3$ for $z=-1$ and $4,6,8$ for
$z=+1$. Exactly $s=1$ has the positive target label, so $\sgn F=g$.

For $N\ge2$, put $n=\lfloor N/2\rfloor$. Use $n$ copies of each
first-layer feature and zero in any remaining neuron. In each of the
$N$ second-layer neurons, assign coefficients $1/n$ to the first
feature group and $-r/n$ to the second. Every second-layer activation
is then $\tanh(F(s))$.
The squared affine coefficient norm of $s$ is $41/2$.
The first-layer energy is at most $82d\lambda^2=1$.
Since $N/n\le3$ and $r\le2$, the second-layer energy is at most
$(N/n)(1+r^2)\le15$.
A later layer with all entries $a/N$ maps a common activation $h$ to
$\tanh(ah)$ and has Frobenius energy $a^2$. A readout with all entries
$1/N$ has weighted energy one and returns that common activation.

We need a margin estimate to choose the remaining layers.
For $d\ge7$, $8\lambda<1/2$. On $[0,1/2]$ we have
$\tanh'(u)>1/2$ and $\tanh u\ge u/2$.
For $s=1$, these imply $a\ge\lambda/2$,
$b-a\ge\lambda/2$, and $b+a\ge3\lambda/2$.
Equation~\eqref{eq:tanh-two}, with denominator at most four, gives
$F(1)\ge3\lambda^3/16$. For $s\in\{3,4,6,8\}$, use
$a\ge3\lambda/2$, $a-b\ge\lambda/2$, and $a+b\ge5\lambda/2$
to get $|F(s)|\ge15\lambda^3/16$. Oddness handles $s=-1$.
Every second-layer activation therefore has absolute value at least
\[
 h_0=\tanh(3\lambda^3/16)\ge3\lambda^3/32.
\]

Any $J\le d$ subsequent gain-one tanh layers retain at least $h_0/2$.
Indeed, $\tanh u\ge u-u^3/3$ for $u\ge0$. Iterating from $h_0$,
all virtual values are at most $h_0$, so the total loss is at most
$Jh_0^3/3<h_0/2$, using $h_0\le3\lambda^3/16$ and $J\le d$. Monotonicity transfers this bound to larger actual
activation magnitudes. For $0<h\le1/2$, concavity and
$\tanh1>3/4$ give $\tanh(2h)>(3/2)h$.
Once $h\ge1/2$, one more gain-two layer makes it strictly greater than
$1/2$ and subsequent such layers preserve that property. Hence
\[
 m=\left\lceil\log_{3/2}\!\left[\frac{32}3(82d)^{3/2}\right]
       \right\rceil+1
\]
final gain-two layers suffice after any preceding gain-one layers.
For $d\ge257$, elementary logarithm estimates give
$m\le25+4\log d<d-2=L-3$, so the available depth suffices.
Use gain one in the other $L-3-m$ layers. The total energy is at most
\[
 1+15+(L-3-m)+4m+1
 =L+14+3m\le L+89+12\log d<2L.
\]
All physical weights are bounded by two.

For the scalar lower bound, absorb the scalar downstream signs into the
first affine vector, using oddness of tanh. Let $\rho>0$ be the least
absolute first preactivation. Dividing its affine vector by $\rho$
gives a unit-margin representation of $\sgn A_t$. Its squared norm is
at least $Q_t$, since this sign function fits the history in
Lemma~\ref{lem:selector}. Since $|\tanh u|\le|u|$, the product of
$\rho$ and all remaining scalar weight magnitudes must be at least
$1/2$. Weighted arithmetic--geometric mean yields
\[
 \kappa_{d,1}\ge L(dQ_t/4)^{1/L}.
\]
But $Q_t\ge(2/3)(4^k-1)$ and
$d(1-4^{-k})>24$ for the present $k$, so
$dQ_t/4\ge4^{k+1}=2^L$. The lower bound is at least $2L$.
\end{proof}

The target comparator itself still has a bounded-weight tanh teacher:
start with $A_M/2^k$, apply tanh, then use $2k$ further scalar hidden
weights two and readout one. The smallest first activation is at least
$(3/4)2^{-k}$. Until it exceeds $1/2$, each following layer multiplies
it by more than $3/2$; after $2k$ layers the resulting lower estimate is
$(3/4)(9/8)^k>1/2$. All coefficients are at most two. The carry
inequalities exclude complexity one exactly as in the near-identity
teacher argument. Thus $C_L(f)=2$ at this depth.

Equal replication of a scalar tanh chain to any width preserves its
energy: duplicate the first row, divide every subsequent scalar
connection and readout by the width, and use the coefficients in
\eqref{eq:energy}. Proposition~\ref{prop:tanh-shortcut} therefore shows
that these replicated scalar implementations are not minimizers in the
larger architecture. It does \emph{not} lower-bound every wider
implementation of $\sgn A_t$, and does not identify the wider posterior's
prediction. Some different wider implementation of a wrong function
could still have lower energy than the displayed competitor.
